% 原像判据示意：输出R中生成集合(a,+inf)与(b,+inf)（a<b），输入X中原像嵌套；(a,b]的原像是两原像之差（蓝色环形区域）。
\begin{tikzpicture}[x=1cm,y=1cm,>=Stealth,
  note/.style={font=\FigureNoteFont\fontsize{7.5}{9}\selectfont,text=FigureMuted},
  tl/.style={font=\fontsize{7.5}{9}\selectfont}]
\draw[FigureStroke,line width=.7pt,fill=white] (0,0) ellipse (1.55 and 1.05);
\node at (-1.3,0.95) {$X$};
\draw[FigureBlue,line width=.7pt,fill=FigureBlueFill] plot[smooth cycle,tension=.8] coordinates {(-0.8,-0.5) (0.2,-0.72) (1.15,-0.3) (1.15,0.45) (0.3,0.72) (-0.75,0.4)};
\draw[FigureGreen,line width=.7pt,fill=FigureGreenFill] plot[smooth cycle,tension=.8] coordinates {(0.2,-0.3) (0.8,-0.25) (0.9,0.25) (0.45,0.42) (0.12,0.1)};
\draw[FigureStroke,line width=.5pt] (0.55,0.15)--(0.9,1.35) node[above,tl,text=FigureGreen!60!black] {$f^{-1}((b,+\infty))$};
\draw[FigureStroke,line width=.5pt] (-0.45,-0.1)--(-0.9,-1.35) node[below,tl,text=FigureBlue] {$f^{-1}((a,b])$};
\begin{scope}[xshift=3.75cm]
  \draw[FigureStroke,line width=.7pt,->] (0,0)--(4.6,0) node[right,tl] {$\mathbb R$};
  \foreach \x/\l in {1.2/a,2.8/b}{\draw[FigureStroke,line width=.6pt] (\x,-.08)--(\x,.08);\node[tl,below] at (\x,-.08) {$\l$};}
  \draw[FigureBlue,line width=2.6pt] (1.2,0.32)--(4.5,0.32);
  \node[note,anchor=south] at (2.0,0.42) {$(a,+\infty)$};
  \draw[FigureGreen,line width=2.6pt] (2.8,0.72)--(4.5,0.72);
  \node[note,anchor=south] at (3.65,0.82) {$(b,+\infty)$};
  \draw[FigureBlue,line width=.6pt,decorate,decoration={brace,mirror,amplitude=3pt}] (1.22,-0.45)--(2.78,-0.45);
  \node[note] at (2.0,-0.78) {$(a,b]$：二者之差};
\end{scope}
\draw[FigureStroke,line width=.7pt,->] (1.7,0.3) to[bend left=18] node[above,tl] {$f$} (3.55,0.3);
\end{tikzpicture}
